% =====================================================================
%  CHAPTER 1: 绪论、儿童少年生长发育规律（对应大纲第一章）
%  内容来源：往届笔记第一章+第二章，参考PPT第一章+第二章
% =====================================================================
\chapter{绪论、儿童少年生长发育规律}
\setcounter{chapter}{1}
\chapterintro{
\textbf{掌握：}课程的特点，儿童少年生长发育的规律及其卫生学意义。\newline
\textbf{熟悉：}课程的性质、任务、研究内容。\newline
\textbf{其他：}课程的研究方法。
}

% ===== 第一部分：课程概述 =====
\section{儿童少年卫生学的性质、任务与研究内容}

\subsection{课程性质与研究对象}

\begin{defbox}
\textbf{儿童少年卫生学（child and adolescent health）：}研究维护和促进儿童少年健康的一门学科。研究儿童少年身心发育随年龄变化的特征，分析影响生长发育的遗传和环境因素，阐明儿童少年机体与学习及生活环境间的相互关系，制定相应卫生要求和措施，预防疾病、增强体质，促进身心健康发育，并为成年期健康奠定良好基础，从而达到提高生命质量的目的。

\textbf{研究对象：}0-24岁儿童少年；重点人群：中小学生群体。
\end{defbox}

\subsection{儿童少年年龄范围的界定}

\begin{keybox}
\textbf{儿童少年年龄分期：}

\begin{table}[H]
\centering
\caption{儿童少年年龄分期}
\begin{tabularx}{\textwidth}{l>{\raggedright\arraybackslash}p{3.5cm}X}
\hline
\textbf{分期} & \textbf{英文名称} & \textbf{年龄范围} \\
\hline
\imp{胎儿期} & fetal period & 从受精卵形成开始到孕40周胎儿娩出 \\
\imp{婴儿期} & infant period & 出生后1年 \\
\imp{新生儿期} & neonatal period & 出生后28天（属于婴儿期） \\
\imp{幼儿期} & toddle period & 出生后第2-3年 \\
\imp{学龄前期} & preschool period & 3-6岁 \\
\imp{学龄期} & school-age period & 6-11/12岁，相当于小学阶段 \\
\imp{青春期} & adolescence & 从青春发动开始到生长基本成熟（10-19岁） \\
\imp{青年期} & youth & 15-24岁 \\
\hline
\end{tabularx}
\end{table}
\end{keybox}

\subsection{儿童少年生物和社会特征}

\begin{keybox}
\textbf{儿童少年三大生物和社会特征：}

\begin{enumerate}[leftmargin=*]
    \item \imp{处于生长发育过程中}——生长发育是儿童少年的基本特征之一，也是社会发展的一面镜子。了解儿童少年生长发育规律，并深入研究其影响因素，是制定儿童少年保健工作的相应政策和行动纲领的重要前提。

    \item \imp{接受学校教育的塑造}——接受教育是儿童少年的一个重要社会特征，学校集体生活构成其重要的生活内容。创造良好的学校物质环境和社会支持环境，建立健康促进学校，开展综合性学校卫生服务，是促进其身心健康和生长发育潜能的重要条件。

    \item \imp{具有社会脆弱性和健康易损性}——在全人口中，儿童少年由于其身心以及各种能力发育的不完善，始终是一个弱势的群体。儿童神经系统发育不成熟，与成人相比，更容易受到有害生物和环境因素影响，从而造成发育和行为异常。
\end{enumerate}
\end{keybox}

\subsection{儿童少年生存与发展权利}

\begin{defbox}
\textbf{儿童少年生存与发展权利：}

\begin{enumerate}[leftmargin=*]
    \item \imp{生存权}：每一个儿童都应该享有的对自己生命权以及维持基本健康生活需求的生活条件保障权。
    \item \imp{发展权}：保障儿童成长过程中的各种需要得到满足，包括儿童有接受一切形式的教育（正规的和非正规的教育）的权利，能够给予儿童的身体、心理、精神、道德与社交发展的相应生活水平。发展权包括参与权、受保护权。
\end{enumerate}
\end{defbox}

\subsection{研究内容与三级预防理论}

\begin{defbox}
\textbf{儿童少年卫生学的四大研究内容：}

\begin{enumerate}[leftmargin=*]
    \item \imp{儿童少年生长发育规律}——生长发育作为人类最为复杂的生物学现象，是本学科研究的永恒主题，也是学科首先要研究和认识的理论问题。包括身体发育和心理发育两个方面。
    \item \imp{儿童少年健康状况及其决定因素}——以保护、促进、增强儿童少年身心健康为宗旨，提出相应的卫生要求和适宜的卫生措施，维护儿童少年的健康，减少疾病，预防死亡，提高生存质量。
    \item \imp{儿童少年卫生服务}——包括：生长发育和健康监测；健康教育和健康促进；常见病防治与健康管理；心理卫生服务；学校营养服务；校园安全管理与伤害及暴力预防；体育与体力活动促进。
    \item \imp{儿童少年卫生学的技术和方法}——人体测量、健康相关体能测试、行为评定等。
\end{enumerate}
\end{defbox}

\begin{keybox}
\textbf{三级预防理论：}

三级预防理论是学科发展的重要理论构架。三级预防是儿童少年健康促进的首要和有效手段，是现代医学为人们提供的健康保障。

\begin{enumerate}[leftmargin=*]
    \item \imp{一级预防（primary prevention）/ 病因预防}：疾病尚未发生时针对致病因素或危险因素采取措施，是预防和消灭疾病的根本措施。是最积极最有效的预防措施。

    \item \imp{二级预防（secondary prevention）/ "三早"预防}：早发现、早诊断、早治疗，是防止或减缓疾病发展采取的措施。

    \item \imp{三级预防（tertiary prevention）/ 临床预防}：防治伤残和促进功能恢复，提高生命质量，延长寿命，降低病死率，主要是对症治疗和康复治疗措施。
\end{enumerate}
\end{keybox}

% ===== 第二部分：研究方法 =====
\section{课程的研究方法}

\begin{tipbox}
\textbf{儿童少年卫生学的研究方法}主要包括：

\begin{enumerate}[leftmargin=*]
    \item \textbf{人体测量方法}——对人体有关部位长度、宽度、厚度和围度的测量，是评价儿童身体匀称度、大小和成分的最简便指标。
    \item \textbf{健康相关体能测试}——采用医学手段检测机体在日常体力活动或运动中的功能。
    \item \textbf{行为评定方法}——通过观察、检查、量表等方法对个体行为问题作出评价。
    \item \textbf{问卷调查与流行病学方法}——横断面设计、监测设计、纵向设计、序列设计等。
\end{enumerate}
\end{tipbox}

% ===== 第三部分：生长发育规律 =====
\section{儿童少年生长发育的规律及其卫生学意义}\difficulty

\begin{keybox}
\textbf{儿童少年生长发育的规律：}

生长发育是儿童少年的基本特征之一，也是社会发展的一面镜子。儿童少年身心发育随年龄变化的特征，受遗传和环境因素共同影响。生长发育是累积和连续的过程，个体从"胎儿—儿童—少年—青壮年—老年"构成完整的生命历程，每一个阶段健康既是前一阶段的延续，也是下一阶段的基础。

\textbf{卫生学意义：}

\begin{enumerate}[leftmargin=*]
    \item 了解儿童少年生长发育规律，并深入研究其影响因素，是制定儿童少年保健工作的相应政策和行动纲领的重要前提。
    \item 儿童少年发展潜力的发挥必须通过整个生命历程中的支持性环境来实现。"支持性环境"是能够激发儿童少年最佳发展潜力的系统、力量和因素，是由贯穿整个生命历程的卫生服务、政策、计划和社区共同参与。
    \item 生长发育规律为制定相应卫生要求和措施、预防疾病、增强体质、促进身心健康发育提供科学依据，并为成年期健康奠定良好基础，从而达到提高生命质量的目的。
\end{enumerate}
\end{keybox}

\section{生长发育的一般规律}

\subsection{生长发育阶段性与连续性的统一}

\begin{defbox}
\textbf{1. 生长发育的阶段性}

\textbf{2. 生长发育的连续性}

（1）\textbf{生长轨迹现象（trajectory of growth）/ 生长管道现象（growth canalization）：}群体儿童少年在正常环境下，生长过程将按遗传潜能决定的方向、速度和目标发育。个体儿童的生长轨迹既与遗传有关，还与疾病及其治疗、营养、体力活动、情感环境等多种因素有关。

（2）\textbf{追赶性生长（catch-up growth）：}处于生长发育过程中的个体受疾病、营养、心理应激等因素作用下，会出现暂时性生长发育的连续性被打破现象，如生长迟缓，一旦这些影响因素解除，机体表现为向原有正常轨道靠近并具有生长发育强烈的倾向，年龄越小越明显，生长加速幅度越大。这种在阻碍生长发育因素解除后出现的生长加速现象。
\end{defbox}

\subsection{生长发育程序性和时间性的协调}

\begin{keybox}
\textbf{1. 生长发育的程序性}

（1）\textbf{运动能力发育的程序性}

①\textbf{头尾发展律（principal of cephalocaudal development）：}婴幼儿时期，粗大动作按抬头、翻身、坐、爬、站、走、跑、跳的发育程序进行。

②\textbf{近侧发展律（principal of proximodistal development）：}粗大动作和精细动作遵循近躯干的四肢肌肉先发育，手的精细动作后发育。

（2）\textbf{线性生长规律：}青春期前的身体线性生长也遵循头尾发展律，2个月龄的胎儿头与躯干比例1:1，出生时1:4，成人1:8。

\textbf{2. 生长发育的个体差异性}
\end{keybox}

\subsection{生长发育不同步性与多样性的平衡}

\begin{keybox}
\textbf{1. 不同组织器官系统的生长不同步性}

（1）\textbf{斯卡蒙生长模式}

①\textbf{一般型：}婴幼儿期增长快；学龄前、学龄期增长平稳；青春期增长又加快；青春后期生长逐步停止。

②\textbf{淋巴系统型：}淋巴结、胸腺、扁桃体等淋巴器官儿童期生长很快，青春期达高峰，以后逐渐衰退，成年时仅相当于高峰时的一半。

③\textbf{神经系统型：}中枢神经系统，视听觉器官、颅骨等仅有一个生长突增期：出生前至学龄前期生长迅速，6岁左右成熟度达90\%左右。

④\textbf{生殖系统型：}除子宫外的生殖器官，青春期前几乎呈停滞状态，青春期一开始迅速加快。

（2）\textbf{子宫生长模式：}子宫、肾上腺出生时较大，其后迅速变小，青春期开始前恢复到出生时的大小，其后迅速增大。

\textbf{2. 生长发育多样性}
\end{keybox}

\subsection{生长发育的高度可塑性}

\begin{defbox}
\textbf{生长发育可塑性（plasticity of growth and development）：}人的结构、功能为适应环境（积极/消极的内外环境）变化和生活经历发生改变的能力。

生长发育的高度可塑性：与年龄、环境和干预敏感期关系密切。
\end{defbox}

% ===== 生长发育概念与指标体系 =====
\section{生长发育概念与指标体系}

\subsection{生长发育基本概念}

\begin{defbox}
\textbf{1. 生长（growth）：}身体各部分及全身在大小、长短和重量上的增加及身体化学成分的变化，含形态生长和化学生长。

\textbf{2. 发育（development）：}身体组织、器官、系统在功能上不断分化与完善过程，含"身"（体格、体力）和"心"（心理、情绪、行为）两个密不可分的方面。

\textbf{3. 成熟（maturity）：}生长和发育达到一个相对完备的阶段，标志个体在形态、生理功能、心理素质等方面都已达到成人水平，具备独立生活和生殖养育下一代的能力。

\textbf{4. 生长发育可塑性（plasticity of growth and development）：}人的结构、功能为适应环境（积极/消极的内外环境）变化和生活经历发生改变的能力。
\end{defbox}

\subsection{生长发育指标体系}

\begin{keybox}
\textbf{1. 体格发育（physical growth）：}身体外部形态的发育，是人体整体发育的重要方面。
\begin{enumerate}[leftmargin=*]
    \item 纵向测量指标
    \item 横向测量指标：围度和径长
    \item 重量测量指标：体重
    \item 体格发育派生指标：BMI、腰高比、腰臀比
\end{enumerate}

\textbf{2. 体能发育指标}
\begin{enumerate}[leftmargin=*]
    \item 生理功能指标：心血管功能指标（心率、脉搏、动态血压）、肺功能指标（呼吸频率、肺活量、最大通气量、最大摄氧量）、肌肉发育指标
    \item 运动能力指标：力量指标、耐力指标、速度指标、灵敏性指标、柔韧性指标
    \item 体能发育派生指标
\end{enumerate}

\textbf{3. 心理行为发育指标：}认知能力指标（感知能力、记忆能力、注意能力、思维能力、执行能力）、情绪状态指标（焦虑、抑郁、恐惧、偏执）、个体发育指标、社会适应能力指标。
\end{keybox}

% ----- 复习题 -----
\reviewcontent{
\section{复习题}

\begin{enumerate}[leftmargin=*, itemsep=8pt]
    \item \textbf{【名词解释】}儿童少年卫生学；生长；发育；成熟；生长发育可塑性；生长轨迹现象；追赶性生长；头尾发展律；近侧发展律；一级预防；二级预防；三级预防
    \item \textbf{【简答题】}简述儿童少年的年龄分期及各期年龄范围。
    \item \textbf{【简答题】}简述儿童少年的三大生物和社会特征。
    \item \textbf{【简答题】}简述儿童少年卫生学的四大研究内容。
    \item \textbf{【简答题】}简述三级预防策略的主要内容。
    \item \textbf{【简答题】}简述斯卡蒙生长模式的四种类型及其特点。
    \item \textbf{【简答题】}试述儿童少年生长发育的规律及其卫生学意义。
\end{enumerate}

\subsection*{参考答案要点}

\begin{enumerate}[leftmargin=*, itemsep=5pt]
    \item \textbf{儿童少年卫生学}：研究维护和促进儿童少年健康的一门学科，研究儿童少年身心发育随年龄变化的特征，分析影响生长发育的遗传和环境因素，阐明儿童少年机体与学习及生活环境间的相互关系，制定相应卫生要求和措施，预防疾病、增强体质，促进身心健康发育，并为成年期健康奠定良好基础，从而达到提高生命质量的目的。研究对象为0-24岁儿童少年，重点人群为中小学生群体。\textbf{生长}：身体各部分及全身在大小、长短和重量上的增加及身体化学成分的变化，含形态生长和化学生长。\textbf{发育}：身体组织、器官、系统在功能上不断分化与完善过程，含"身"（体格、体力）和"心"（心理、情绪、行为）两个密不可分的方面。\textbf{成熟}：生长和发育达到一个相对完备的阶段，标志个体在形态、生理功能、心理素质等方面都已达到成人水平，具备独立生活和生殖养育下一代的能力。\textbf{生长发育可塑性}：人的结构、功能为适应环境（积极/消极的内外环境）变化和生活经历发生改变的能力。\textbf{生长轨迹现象}：群体儿童少年在正常环境下，生长过程将按遗传潜能决定的方向、速度和目标发育。\textbf{追赶性生长}：处于生长发育过程中的个体受疾病、营养、心理应激等因素作用下，出现暂时性生长发育的连续性被打破现象，一旦这些影响因素解除，机体表现为向原有正常轨道靠近并具有生长发育强烈的倾向，年龄越小越明显，生长加速幅度越大。\textbf{头尾发展律}：婴幼儿时期，粗大动作按抬头、翻身、坐、爬、站、走、跑、跳的发育程序进行。\textbf{近侧发展律}：粗大动作和精细动作遵循近躯干的四肢肌肉先发育，手的精细动作后发育。\textbf{一级预防（primary prevention）/ 病因预防}：疾病尚未发生时针对致病因素或危险因素采取措施，是预防和消灭疾病的根本措施，是最积极最有效的预防措施。\textbf{二级预防（secondary prevention）/ "三早"预防}：早发现、早诊断、早治疗，是防止或减缓疾病发展采取的措施。\textbf{三级预防（tertiary prevention）/ 临床预防}：防治伤残和促进功能恢复，提高生命质量，延长寿命，降低病死率，主要是对症治疗和康复治疗措施。

    \item 八个分期：胎儿期（受精卵形成至孕40周娩出）、婴儿期（出生后1年）、新生儿期（出生后28天，属婴儿期）、幼儿期（出生后第2-3年）、学龄前期（3-6岁）、学龄期（6-11/12岁，相当于小学阶段）、青春期（10-19岁，从青春发动开始到生长基本成熟）、青年期（15-24岁）。

    \item (1) 处于生长发育过程中；(2) 接受学校教育的塑造；(3) 具有社会脆弱性和健康易损性。

    \item (1) 儿童少年生长发育规律；(2) 儿童少年健康状况及其决定因素；(3) 儿童少年卫生服务；(4) 儿童少年卫生学的技术和方法。

    \item 一级预防（病因预防）：疾病尚未发生时针对致病因素或危险因素采取措施；二级预防（"三早"预防）：早发现、早诊断、早治疗；三级预防（临床预防）：防治伤残和促进功能恢复。

    \item ①一般型：婴幼儿期增长快，学龄期平稳，青春期加快，后期停止；②淋巴系统型：儿童期生长快，青春期达高峰，以后衰退；③神经系统型：仅有一个生长突增期，6岁左右成熟度达90\%；④生殖系统型：青春期前几乎停滞，青春期开始迅速加快。子宫生长模式：出生时较大，迅速变小，青春期开始前恢复，其后迅速增大。

    \item 生长发育是累积和连续的过程，受遗传和环境因素共同影响。卫生学意义：(1)了解规律并研究影响因素，是制定政策和行动纲领的重要前提；(2)发展潜力的发挥需通过支持性环境实现；(3)为制定卫生要求和措施、预防疾病、增强体质、促进身心健康发育提供科学依据，为成年期健康奠定基础。
\end{enumerate}
}

\newpage
